\section{The next proof obligation}

The finite step now returns the same analytic radius and the
same positive connected-carrier exponent. Its amplitude estimate
still has a coefficient larger than one multiplying the inherited
activity norm. The background-dependent determinant and delta
terms are also not small merely because the coupling is small.
These statements locate the remaining small-field induction
problem: extract the vacuum and allowed marginal terms, establish
the coupling recursion, and prove a uniform bound on the
remaining interactions through arbitrarily many steps.

The high-field-degree estimates provide one part of that
recursion: coefficientwise fifth-order background remainders
contract at a sufficiently large fixed macro ratio, and
the actual quantum step has a contractive response to
differences with matching low jets. Its Taylor subtraction
is bounded without a radius loss. This does not give an
invariant high-degree input class, because the fluctuation
integral can regenerate lower-degree terms.
The low-degree pieces still include kernels on arbitrarily
many carrier geometries, and the separate pieces need not be
gauge invariant. A complete extraction must combine the
field-degree estimate with spatial locality and the nonlinear
gauge identities, and include the generated low jets in the
recurrence.

The total correction has the local gauge identity needed for
that analysis. A scalar marginal gauge-action coefficient still
requires the appropriate lattice symmetry classification; a
Ward identity alone does not supply it. The large-field
complements, mixed-history matching, observable estimates,
continuum nontriviality and reconstruction remain separate
required parts of the four-dimensional existence theorem.

\begin{thebibliography}{9}
\bibitem{priorcd}
\emph{Toward a four-dimensional Yang--Mills construction},
research Notebooks 15--27 in this workspace.
The homogeneous classical support families, strict overlapping
Gaussian expansion, gauge identities, supported pre-gauge
representative, auxiliary grouping and exact composite macrostep
are the principal inputs.
\bibitem{balaban116cd}
T.~Ba\l aban,
\emph{Renormalization group approach to lattice gauge field
theories. II. Cluster expansions},
Commun. Math. Phys. \textbf{116} (1988), 1--22.
\href{https://scispace.com/pdf/renormalization-group-approach-to-lattice-gauge-field-4penm4k7hc.pdf}{Original article copy}.
\bibitem{sourcecd}
E.~Sussman and collaborators,
\href{https://github.com/ethanwsussman/Yang--Mills-LLM-experiment}{Yang--Mills-LLM-experiment},
deposited raw accounts at commit
\texttt{e392ef45013a581664bab76b90867170b30fe10e}.
\end{thebibliography}
\end{document}
